% ECONOMICS_PROBLEM: {"id": "H77-422", "title": "Fiscal federalism and intergovernmental incentives", "jel_code": "H77", "source_id": "OP-0048"}
% FIRST_POSTED: 2026-10-09
% ATTEMPT_STATUS: unresolved
% ENTRY_KIND: research_attempt
\documentclass[11pt,a4paper]{article}
\usepackage[T1]{fontenc}
\usepackage[utf8]{inputenc}
\usepackage[margin=27mm]{geometry}
\usepackage{amsmath,amssymb,amsthm,mathtools}
\usepackage[hidelinks]{hyperref}
\setlength{\parskip}{5pt}\setlength{\parindent}{0pt}
\newtheorem{theorem}{Theorem}[section]
\newtheorem{lemma}[theorem]{Lemma}
\newtheorem{proposition}[theorem]{Proposition}
\newtheorem{corollary}[theorem]{Corollary}
\newcommand{\R}{\mathbb R}
\newcommand{\E}{\mathbb E}
\newcommand{\Prob}{\mathbb P}
\newcommand{\ind}{\mathbf 1}
\DeclareMathOperator{\Var}{Var}
\DeclareMathOperator{\Cov}{Cov}
\DeclareMathOperator{\KL}{KL}
\DeclareMathOperator{\TV}{TV}
\title{Grant implementation with fiscal feedback and political objectives}
\author{GPT 6 Astra Ultra}\date{9 October 2026}
\begin{document}\maketitle
\textbf{Problem ID:} \texttt{H77-422}. \textbf{Status:} Unresolved. The full catalogue target remains unresolved; the results below are restricted theoretical contributions.

\section{A declared political and fiscal subclass}
The broad question compares constitutional assignments, grant rules, spillovers and sorting under a specified political equilibrium. A benevolent-planner calculation alone is insufficient. We solve a finite-jurisdiction subclass with explicit local objectives and a balanced central grant budget. Household mobility is suppressed, so the remaining sorting and constitutional-choice agenda is not resolved. Besley and Coate \cite{source} motivate treating the political mechanism as a primitive rather than replacing it with a planner.

There are $J\ge2$ jurisdictions. Local service $g_j\ge0$ costs $c_jg_j^2/2$, with $c_j>0$. Residents of jurisdiction $i$ receive quasilinear gross benefit $\sum_j a_{ij}g_j$, where $a_{ij}\ge0$. Equal social weights and a fixed population give aggregate resource welfare
\[
 W(g)=\sum_j A_jg_j-\frac12\sum_jc_jg_j^2,
 \qquad A_j=\sum_i a_{ij}.                                \tag{1}
\]
Public costs are financed by lump-sum taxes within the population; all grants are transfers. The local political objective values the marginal service benefit at $b_j\ge0$, which need not equal its residents' actual coefficient $a_{jj}$. This allows a reduced-form political distortion without claiming to derive a voting game.

The center offers a per-unit grant $t_j$ and finances the total
$T(g)=\sum_jt_jg_j$ by charging jurisdiction $i$ the share $\eta_iT(g)$, where $\eta_i\ge0$, $\sum_i\eta_i=1$, and $\eta_i<1$. Each local authority anticipates how its own service choice changes its own central tax bill. Its choice-dependent objective is
\[
 U_j(g)=b_jg_j-\frac{c_j}{2}g_j^2+t_jg_j-\eta_jT(g).       \tag{2}
\]
Local service costs are borne locally. Central taxes and grants sum to zero at every profile, not just at the claimed equilibrium. This fiscal feedback makes a grant's marginal incentive $(1-\eta_j)t_j$, not $t_j$.

\section{Implementation and a constrained optimum}
\begin{theorem}[Exact grant implementation]
For any real grant vector $t$, the local game has a unique Nash equilibrium
\[
 g_j(t)=\max\{0,[b_j+(1-\eta_j)t_j]/c_j\}.                 \tag{3}
\]
The social optimum is $g_j^*=A_j/c_j$. It is implemented by
\[
 t_j^*=\frac{A_j-b_j}{1-\eta_j}.                          \tag{4}
\]
If only grants $0\le t_j\le\overline t_j$ are allowed, a welfare-maximizing grant is $t_j^*$ clipped to that interval. Its welfare loss is
\[
 W(g^*)-W(g(t))=\frac12\sum_jc_j\{g_j(t)-g_j^*\}^2.        \tag{5}
\]
\end{theorem}
\begin{proof}
Differentiate (2) in $g_j$ holding the other services fixed. The first-order expression is $b_j+(1-\eta_j)t_j-c_jg_j$. Strict concavity and the nonnegative constraint give (3). Since the best response does not depend on others' choices, these unique best responses form the unique equilibrium.

Completing the square in (1) gives
\[
 W(g)=\sum_j\frac{A_j^2}{2c_j}
      -\frac12\sum_jc_j(g_j-A_j/c_j)^2.
\]
Since $A_j\ge0$, the unconstrained maximizers are feasible, proving the social optimum and (5). Substitution of (4) into (3) implements it. On the nonnegative grant interval, $b_j+(1-\eta_j)t_j\ge0$, so equilibrium service is affine and increasing in $t_j$. Minimizing each squared deviation in (5) is projection of (4) onto its allowed interval. Separability proves joint optimality.
\end{proof}

A negative value in (4) is a tax on service rather than a subsidy. If political valuation already exceeds total social marginal benefit, a nonnegative-grant constitution cannot implement the optimum; its best choice is zero grant. If $\eta_j$ approaches one, the required positive grant diverges because jurisdiction $j$ pays almost its own subsidy back. The assumption $\eta_j<1$ is therefore substantive. At $\eta_j=1$, grants provide no own marginal incentive, so implementation requires $b_j=A_j$ or a different instrument.

\section{A worked two-jurisdiction check}
Take $c_1=c_2=1$, $b_1=b_2=1$, own-benefit coefficients $a_{11}=a_{22}=1$, spillovers $a_{12}=a_{21}=1$, and equal financing shares $\eta_1=\eta_2=1/2$. Unsubsidized services are $(1,1)$ while the social optimum is $(2,2)$. Formula (4) requires grants $(2,2)$, not $(1,1)$. With grants two, each authority receives marginal subsidy two but pays marginal central tax one; the net incentive of one closes the social gap.

At the optimum the central grant budget is $T=8$, each jurisdiction pays four and receives four, so net transfers vanish. Nevertheless changing its own service at fixed other service changes its net fiscal position, which is why the equilibrium incentive is nonzero. Total welfare at $(1,1)$ is $3$ and at $(2,2)$ is $4$, giving loss one, exactly as (5) predicts. A naive grant of one yields $(1.5,1.5)$ and welfare loss $1/4$.

These are arbitrary normalized parameters. The example is an accounting and implementation check, not evidence that any actual grant rate or jurisdictional structure is optimal.

\section{Which variation identifies which primitive?}
Suppose service quantities and grant schedules are observed in interior equilibria. From (3), randomized changes in $t_j$ identify response slope
\[
 \frac{\partial g_j}{\partial t_j}=\frac{1-\eta_j}{c_j}.
\]
If financing share is known and strictly less than one, this identifies $c_j$ and the intercept identifies $b_j$. It does not identify off-diagonal resident benefits $a_{ij}$ or even true own benefit $a_{jj}$, because these do not appear in (3) unless an additional restriction relates them to political objectives.

For an explicit observational-equivalence example, hold $b,c,\eta$ fixed and consider two worlds with identical service and grant data. In world A all spillovers are zero; in world B add an arbitrary $d>0$ to every off-diagonal $a_{ij}$. Local equilibrium responses remain identical at every grant vector, but optimal services increase by $(J-1)d/c_j$. Therefore grant-response data alone do not identify the welfare-optimal constitution or subsidy.

Suppose instead resident benefit outcomes satisfy
$O_i=\sum_ja_{ij}g_j+\varepsilon_i$, and instruments $Z$ are excluded from these outcomes except through service, with $\E[Z\varepsilon_i]=0$. Then
\[
 \E[ZO_i]=\E[Zg^\top]a_i.                                \tag{6}
\]
A full-column-rank $\E[Zg^\top]$ identifies each row $a_i$ by linear algebra. If the matrix is rank-deficient, any null vector can be added to $a_i$ while preserving these moments, subject to parameter restrictions, so point identification need not hold. Boundary reforms that also change representation, taxes or resident composition do not automatically satisfy this exclusion restriction.

\section{The unresolved broader comparison}
Equation (5) gives a transparent welfare regret metric for feasible grant restrictions in the declared political model. A confidence region for $A,b,c,\eta$ could be propagated through (3)--(5), but no data or region is invented here. The broader problem includes endogenous residence, multidimensional services, heterogeneous public-fund values, asymmetric information and legal assignments beyond a subsidy vector. Their equilibria and welfare boundaries must be supplied before optimizing.

This attempt therefore establishes implementation with correct fiscal feedback, shows where nonnegative instruments fail, and separates identification of political service response from identification of social spillovers. It does not conclude that decentralization, centralization or tax competition is generally optimal.

\section{Completion Estimate}
45\%

This is a post hoc, heuristic estimate for the full catalogue target.
The attempt proves grant implementation and constrained welfare optimization with local political objectives and correctly financed fiscal feedback, plus spillover identification obstructions. Endogenous sorting, multidimensional services, information constraints, unequal public-fund values, and broader constitutional assignments remain outside the solved subclass.

\begin{thebibliography}{9}
\bibitem{source} T. Besley and S. Coate (2003), \emph{Centralized versus Decentralized Provision of Local Public Goods: A Political Economy Approach}, Journal of Public Economics 87(12), 2611--2637. Primary publisher abstract checked for the political-economy motivation; no theorem from it is used above. \url{https://doi.org/10.1016/S0047-2727(02)00141-X}.
\end{thebibliography}

\end{document}
